% Apunts de teoria en PDF · 4t ESO · Unitat 7: Funcions elementals
% Mateix estil que els apunts de Batxillerat (franja i caixes de teoria en blau, exemples en roig), però més visuals,
% curts i amb dibuixos, per a l'alumnat de 3r ESO. Es compila amb pdflatex dues vegades (per als enllaços de l'índex),
% amb logo-fentmates.pdf (_eines/logo/pdf/horitzontal-clar.pdf) al costat.
\documentclass[11pt,a4paper]{article}
\usepackage[utf8]{inputenc}
\usepackage[T1]{fontenc}
\usepackage[catalan]{babel}
\usepackage[sfdefault,book,lining]{FiraSans}
\usepackage[italic,noplusnominus,noequal]{mathastext}
\linespread{1.18}
\usepackage[a4paper,top=2.4cm,bottom=2.2cm,left=1.6cm,right=1.6cm,headheight=1.1cm,headsep=0.5cm,footskip=1cm]{geometry}
\usepackage{amsmath,amssymb,cancel}
\usepackage{icomma}
\usepackage{graphicx}
\usepackage[table]{xcolor}
\usepackage{tikz}\usetikzlibrary{calc,positioning,arrows.meta,decorations.pathreplacing,shapes.geometric}
\usepackage{fontawesome5}
\usepackage{ccicons}
\usepackage{fancyhdr}
\usepackage{enumitem}
\usepackage{array}
\usepackage{needspace}
\usepackage[most]{tcolorbox}
\usepackage{titlesec}
\usepackage[hidelinks]{hyperref}
\hypersetup{pdftitle={Apunts de teoria. Unitat 7: Funcions elementals. 4t ESO},pdfsubject={Matemàtiques A 4t ESO},pdfkeywords={funció lineal, pendent, equació punt-pendent, funció quadràtica, paràbola, vèrtex, funcions radicals, proporcionalitat inversa, hipèrbola, asímptotes, translacions, funció exponencial, problemes}}

% colors de la web (bloc d'Àlgebra: blau)
\definecolor{fmnavy}{HTML}{1B2748}
\definecolor{fmblue}{HTML}{3D6AA0}
\definecolor{fmbluel}{HTML}{4E81BD}
\definecolor{fmyellow}{HTML}{FFC94D}
\definecolor{fmgold}{HTML}{E9A919}
\definecolor{fmaccent}{HTML}{FF7A50}
\definecolor{fmdef}{HTML}{1D4F73}\definecolor{fmdefbg}{HTML}{EEF4F9}
\definecolor{fmex}{HTML}{A3203A}\definecolor{fmexbg}{HTML}{FBF0F2}
\definecolor{fmsf}{HTML}{C2410C}
\definecolor{fmok}{HTML}{15803D}
\definecolor{fmmuted}{HTML}{5D6685}
\definecolor{fmsurface}{HTML}{EEF1F8}
\definecolor{fmA}{HTML}{2F6FD6}\definecolor{fmB}{HTML}{EA6A10}\definecolor{fmR}{HTML}{E11D48}\definecolor{fmK}{HTML}{16A34A}\definecolor{fmV}{HTML}{7C3AED}

% lletres
\newcommand{\comic}{\fontfamily{ComicNeue-TLF}\selectfont}
\newcommand{\heavy}{\fontseries{eb}\selectfont}
\newcommand{\semi}{\fontseries{sb}\selectfont}
\DeclareUnicodeCharacter{2212}{-}

\setlength{\parindent}{0pt}
\setlength{\parskip}{7pt}
\setlength{\lineskip}{4pt}\setlength{\lineskiplimit}{4pt}
\setlength{\emergencystretch}{3em}
\setlist{itemsep=3pt,topsep=3pt,leftmargin=1.4em}
\setcounter{secnumdepth}{1}
\renewcommand{\thesection}{7.\arabic{section}}
\titleformat{\section}[hang]{\color{fmnavy}\fontsize{23}{27}\heavy\selectfont}{\tcbox[on line,colback=fmblue,colframe=fmblue,arc=3mm,boxsep=2pt,left=6pt,right=6pt,top=3pt,bottom=3pt]{\color{white}\thesection}}{0.45em}{}[\vspace{5pt}{\color{fmblue}\titlerule[2pt]}]
\titlespacing*{\section}{0pt}{4pt}{6pt}
\titleformat{\subsection}{\color{fmblue}\Large\semi}{}{0pt}{\faCaretRight\ }
\titlespacing*{\subsection}{0pt}{8pt}{2pt}

% frase de cada apartat, «a mà»
\newcommand{\frase}[1]{\par\vspace{-2pt}{\comic\Large\color{fmbluel}\faCommentDots[regular]\ \ #1\par}\vspace{2pt}}
\newcommand{\seccio}[3]{\Needspace{16\baselineskip}\section{#2}\hypertarget{#1}{}\frase{#3}}
% notes a mà dins dels exemples
\newcommand{\nota}[1]{{\comic\small\color{fmex!85!black}#1}}
\newcommand{\notab}[1]{{\comic\small\color{fmblue}#1}}
% bé / malament
\newcommand{\be}[1]{{\color{fmok}\faCheckCircle}\ \ensuremath{#1}}
\newcommand{\mal}[1]{{\color{fmR}\faTimesCircle}\ {\color{fmR}\ensuremath{#1}}}
% pas numerat
\newcommand{\pas}[1]{\tikz[baseline=(p.base)]\node[circle,fill=fmblue,text=white,inner sep=1.2pt,minimum size=1.25em,font=\small\heavy](p){#1};}
% etiqueta de color
\newcommand{\chip}[2]{\tcbox[on line,colback=#1!14,colframe=#1!60,arc=2mm,boxsep=1pt,left=5pt,right=5pt,top=1pt,bottom=1pt]{#2}}
\newcommand{\fmkey}[1]{\colorbox{fmyellow!55}{\strut #1}}
\newcommand{\cA}[1]{{\color{fmA}#1}}\newcommand{\cB}[1]{{\color{fmB}#1}}\newcommand{\cK}[1]{{\color{fmK}#1}}\newcommand{\cR}[1]{{\color{fmR}#1}}\newcommand{\cV}[1]{{\color{fmV}#1}}

% caixes
\tcbset{fmbase/.style={enhanced,arc=3mm,boxrule=1.2pt,left=4mm,right=4mm,top=2mm,bottom=2.5mm,before skip=8pt,after skip=8pt,fonttitle=\heavy\large,before upper=\raggedright}}
\newtcolorbox{idea}[1][Idea clau]{fmbase,colback=fmdefbg,colframe=fmdef,coltitle=white,title={\faKey\ \ #1}}
\newtcolorbox{exemple}[1][Exemple]{fmbase,colback=fmexbg,colframe=fmex,colbacktitle=fmexbg,coltitle=fmex,title={\faPencil*\ \ #1},titlerule=0.5pt,titlerule style={fmex!40,dashed}}
\newtcolorbox{compte}[1][Compte!]{fmbase,colback=fmsf!6,colframe=fmsf,coltitle=white,fonttitle=\comic\bfseries\Large,title={\faExclamationTriangle\ \ #1}}
\newtcolorbox{truc}[1][Truc]{fmbase,colback=fmyellow!20,colframe=fmgold,colbacktitle=fmyellow,coltitle=fmnavy,fonttitle=\comic\bfseries\Large,title={\faLightbulb[regular]\ \ #1}}

% dibuix de la targeta de la unitat, en blanc
\newcommand{\fmcard}{\begin{tikzpicture}[x=0.062cm,y=-0.062cm,white,line width=3pt,line cap=round]
  \begin{scope}[transparency group,opacity=0.15]
  \node[anchor=base west,inner sep=0pt,font=\fontsize{28}{28}\selectfont] at (14,44) {$y=\sqrt x\quad y=\frac kx$};
  \draw[line width=2pt] (130,105) -- (235,105) (175,115) -- (175,25);
  \draw[line width=3.4pt] (140,30) parabola bend (175,95) (210,30);
  \draw[line width=3.4pt] (180,102) .. controls (200,100) and (215,80) .. (232,30);
  \end{scope}
\end{tikzpicture}}
\newcommand{\fmbrand}{\raisebox{-0.25\height}{\includegraphics[height=0.5cm]{logo-fentmates.pdf}}}

% regla de Ruffini: #1 = a, #2 = coeficients, #3 = fila del mig (sense el primer), #4 = fila de baix, #5 = nombre de coeficients
\newcommand{\ruffini}[5]{\begin{tikzpicture}[x=1.25cm,y=0.85cm,every node/.style={font=\large}]
  \foreach \v [count=\i] in {#2} \node (t\i) at (\i,0) {$\v$};
  \foreach \v [count=\i] in {#3} {\pgfmathtruncatemacro\j{\i+1}\node (m\j) at (\j,-1) {$\v$};}
  \foreach \v [count=\i] in {#4} \node (b\i) at (\i,-2.1) {$\v$};
  \node[text=fmB] at (0,-1) {$#1$};
  \draw[line width=1.1pt,fmnavy] (0.45,0.5)--(0.45,-1.5) (0.45,-1.5)--(#5+0.5,-1.5);
  \node[draw=fmR,line width=1.3pt,rounded corners=2mm,minimum size=0.75cm] at (#5,-2.1) {};
  \draw[-{Latex[length=2.2mm]},fmK,line width=1.1pt] (t1.south) -- (b1.north);
  \pgfmathtruncatemacro\last{#5-1}
  \foreach \i in {1,...,\last} {\pgfmathtruncatemacro\j{\i+1}
    \draw[-{Latex[length=2.2mm]},fmB,line width=1.1pt,shorten >=1pt] (b\i.north east) -- (m\j.south west);
    \draw[-{Latex[length=2.2mm]},fmA!70,line width=1pt,dashed,shorten >=2pt] ([xshift=8pt]t\j.south) -- ([xshift=8pt]b\j.north);}
\end{tikzpicture}}

% capçalera i peu
\pagestyle{fancy}\fancyhf{}
\fancyhead[L]{\small\fmbrand}
\fancyhead[R]{\small\color{fmmuted}4t ESO · Unitat 7: Funcions elementals}
\fancyfoot[L]{\scriptsize\color{fmmuted}\href{https://creativecommons.org/licenses/by-nc-sa/4.0/deed.ca}{\ccbyncsa}\ \ FentMates · fentmates.org · CC BY-NC-SA 4.0}
\fancyfoot[R]{\small\color{fmmuted}\thepage}
\renewcommand{\headrulewidth}{0.4pt}\renewcommand{\headrule}{\color{fmmuted!40}\hrule height \headrulewidth}
\fancypagestyle{primera}{\fancyhf{}\renewcommand{\headrulewidth}{0pt}
  \fancyfoot[L]{\scriptsize\color{fmmuted}\href{https://creativecommons.org/licenses/by-nc-sa/4.0/deed.ca}{\ccbyncsa}\ \ FentMates · fentmates.org · CC BY-NC-SA 4.0}
  \fancyfoot[R]{\small\color{fmmuted}\thepage}}

% targeta de l'índex
\newcommand{\itile}[4]{\begin{tcolorbox}[enhanced,colback=white,colframe=fmblue!35,arc=3mm,boxrule=1pt,boxsep=0pt,left=6pt,right=6pt,top=5pt,bottom=5pt,before skip=0pt,after skip=7pt]
  \hyperlink{#1}{\begin{minipage}[c]{0.9cm}\centering{\color{fmblue}\LARGE #3}\end{minipage}\hspace{6pt}%
  \begin{minipage}[c]{0.78\linewidth}{\color{fmblue}\heavy #2}\\[-2pt]{\color{fmnavy}\semi\large #4}\end{minipage}}\end{tcolorbox}}

% taula amb vores arrodonides: el contingut es retalla amb la forma (sense @{} a les vores ni entre columnes)
\newsavebox\tabclipbox
\newcommand{\taulaclip}[3][3mm]{\sbox\tabclipbox{#3}\begin{tikzpicture}[baseline=(tcn.base)]
  \begin{scope}\clip[rounded corners=#1] (0,-\dp\tabclipbox) rectangle (\wd\tabclipbox,\ht\tabclipbox);
  \node[inner sep=0pt,outer sep=0pt,anchor=base west] (tcn) at (0,0) {\usebox\tabclipbox};\end{scope}
  \draw[#2,rounded corners=#1] (0,-\dp\tabclipbox) rectangle (\wd\tabclipbox,\ht\tabclipbox);
\end{tikzpicture}}
\newcommand{\taularod}[3][fmblue]{\taulaclip{#1,line width=1.2pt}{\renewcommand{\arraystretch}{1.5}\begin{tabular}{#2}#3\end{tabular}}}
% balança
\newcommand{\balanca}[3]{\begin{tikzpicture}[scale=#3,transform shape,every node/.style={anchor=center}]
  \fill[fmnavy] (-0.45,-1.9) rectangle (0.45,-1.75);
  \fill[fmnavy!80] (-0.12,-1.75) rectangle (0.12,0);
  \fill[fmgold] (0,0.05) circle (0.13);
  \draw[fmnavy,line width=2.4pt,line cap=round] (-2,0) -- (2,0);
  \foreach \s in {-1,1} {
    \draw[fmmuted,line width=0.8pt] (\s*2,0) -- (\s*2-0.75,-1.05) (\s*2,0) -- (\s*2+0.75,-1.05);
    \fill[fmyellow!70,draw=fmgold,line width=1pt] (\s*2-0.95,-1.05) arc[start angle=180,end angle=360,x radius=0.95,y radius=0.32] -- cycle;}
  \node[font=\Large,fill=fmA!15,draw=fmA,rounded corners=1.5mm,inner sep=3pt] at (-2,-0.68) {$#1$};
  \node[font=\Large,fill=fmB!15,draw=fmB,rounded corners=1.5mm,inner sep=3pt] at (2,-0.68) {$#2$};
\end{tikzpicture}}
% eixos amb quadrícula (dins d'un tikzpicture): xmin, xmax, ymin, ymax
\newcommand{\eixos}[4]{\draw[fmmuted!22,line width=0.4pt] (#1,#3) grid (#2,#4);
  \draw[-{Latex[length=2mm]},fmmuted,line width=0.8pt] (#1-0.2,0)--(#2+0.4,0) node[right,font=\small]{$x$};
  \draw[-{Latex[length=2mm]},fmmuted,line width=0.8pt] (0,#3-0.2)--(0,#4+0.4) node[above,font=\small]{$y$};}
% taula de traduccions (diccionari) amb les vores arrodonides
\newcommand{\dic}[2]{#1 & {\large$#2$}\\}
\newcommand{\dicratio}{0.62}
\newcommand{\dichead}{\begin{tabular}{@{}>{\raggedright}p{\dimexpr\dicratio\linewidth-12pt}>{\centering\arraybackslash}p{\dimexpr\linewidth-\dicratio\linewidth-12pt}@{}}\faComment[regular]\ Ho diem així\dots & \dots\ així\end{tabular}}
\newenvironment{dicbox}{\begin{tcolorbox}[enhanced,arc=3mm,boxrule=1.2pt,colframe=fmblue,colback=white,coltitle=white,fonttitle=\heavy,
  title={\dichead},left=0pt,right=0pt,top=0pt,bottom=1.5mm,boxsep=0pt,lefttitle=6pt,righttitle=6pt,before skip=0pt,after skip=0pt]\rowcolors{1}{white}{fmsurface}%
  \begin{tabular}{@{\hspace{6pt}}>{\raggedright}p{\dimexpr\dicratio\linewidth-12pt}>{\centering\arraybackslash}p{\dimexpr\linewidth-\dicratio\linewidth-12pt}@{\hspace{6pt}}}}
  {\end{tabular}\end{tcolorbox}}
% resum
\newcommand{\rtile}[3]{\begin{tcolorbox}[enhanced,colback=white,colframe=fmblue!45,arc=3mm,boxrule=1pt,left=3mm,right=3mm,top=2.5mm,bottom=2.5mm,fontupper=\large,
  title={\tcbox[on line,colback=fmblue,colframe=fmblue,arc=2mm,boxsep=0pt,left=4pt,right=4pt,top=1pt,bottom=1pt]{\color{white}\small\heavy #1}\ \ #2},coltitle=fmnavy,colbacktitle=fmdefbg,fonttitle=\semi,toptitle=0.5mm,bottomtitle=0.5mm]#3\end{tcolorbox}}
% portada: #1 subtítol a mà
\newcommand{\portada}[1]{\thispagestyle{primera}
\begin{tikzpicture}[remember picture,overlay]
  \fill[fmblue] (current page.north west) rectangle ([yshift=-9.4cm]current page.north east);
  \begin{scope}
    \clip (current page.north west) rectangle ([yshift=-9.4cm]current page.north east);
    \node[anchor=north west,inner sep=0pt] at ([xshift=6.6cm,yshift=-1.3cm]current page.north west) {\fmcard};
  \end{scope}
  \node[anchor=north west,inner sep=0pt] at ([xshift=1.6cm,yshift=-1.1cm]current page.north west) {\tcbox[colback=white,colframe=white,arc=2.5mm,boxsep=0pt,left=7pt,right=7pt,top=5pt,bottom=5pt]{\includegraphics[height=0.75cm]{logo-fentmates.pdf}}};
\end{tikzpicture}
\vspace*{0.55cm}
\tcbox[colback=white!15!fmblue,colframe=white!15!fmblue,arc=2.5mm,boxsep=0pt,left=6pt,right=6pt,top=3pt,bottom=3pt]{\color{white}\small\bfseries Apunts de teoria · 4t ESO · Matemàtiques A}

\vspace{2pt}
{\raggedright\color{white}\fontsize{30}{36}\heavy\selectfont{\color{fmyellow}Unitat 7}\\[2pt]Funcions elementals\par}
\vspace{8pt}
{\color{white}\comic\Large #1\par}
\vspace*{\dimexpr 7.4cm-\pagetotal\relax}
}
\newcommand{\llegenda}{\par\vspace{16pt}
{\color{fmnavy}\heavy\Large\faPalette\ \ Com es llegeixen aquests apunts}\par\smallskip
\begin{tcbraster}[raster columns=4,raster equal height,raster column skip=8pt,raster before skip=12pt]
\tcbox[colback=fmdefbg,colframe=fmdef,arc=3mm,boxrule=1.2pt,halign=center,width=\linewidth]{{\color{fmdef}\Large\faKey}\\[2pt]{\color{fmdef}\heavy Idea clau}\\[-1pt]{\small el que cal saber}}
\tcbox[colback=fmexbg,colframe=fmex,arc=3mm,boxrule=1.2pt,halign=center,width=\linewidth]{{\color{fmex}\Large\faPencil*}\\[2pt]{\color{fmex}\heavy Exemple}\\[-1pt]{\small pas a pas}}
\tcbox[colback=fmsf!6,colframe=fmsf,arc=3mm,boxrule=1.2pt,halign=center,width=\linewidth]{{\color{fmsf}\Large\faExclamationTriangle}\\[2pt]{\color{fmsf}\heavy Compte!}\\[-1pt]{\small errors típics}}
\tcbox[colback=fmyellow!20,colframe=fmgold,arc=3mm,boxrule=1.2pt,halign=center,width=\linewidth]{{\color{fmgold!80!black}\Large\faLightbulb[regular]}\\[2pt]{\color{fmnavy}\heavy Truc}\\[-1pt]{\small per a recordar-ho}}
\end{tcbraster}}
\newcommand{\resumtitol}{\clearpage
\hypertarget{resum}{}
{\color{fmnavy}\fontsize{23}{27}\heavy\selectfont\faClipboardList\ \ En resum}\par\vspace{5pt}{\color{fmblue}\titlerule[2pt]}
\frase{Tota la unitat en una pàgina.}}
% els 5 passos per a resoldre problemes: #2 = text del pas 2
\newcommand{\passosproblema}[2]{\begin{center}
\begin{tikzpicture}[s/.style={draw=fmblue,fill=fmdefbg,rounded corners=3mm,line width=1.2pt,minimum width=2.75cm,minimum height=2.45cm,align=center,font=\small},
  ic/.style={font=\LARGE,text=fmblue}]
  \foreach \i/\ic/\t in {1/\faBookOpen/{\textbf{Llegim}\\i entenem},2/\faQuestionCircle/{#1},3/\faEquals/{#2},4/\faCogs/{\textbf{Resolem}},5/\faCheckDouble/{\textbf{Comprovem}\\i responem}} {
    \node[s] (s\i) at (3.3*\i,0) {};
    \node[ic] at ([yshift=-0.5cm]s\i.north) {\ic};
    \node[align=center,font=\small,anchor=north] at ([yshift=-0.95cm]s\i.north) {\t};
    \node[circle,fill=fmblue,text=white,font=\small\heavy,inner sep=1.5pt,minimum size=1.3em] at (s\i.north west) {\i};}
  \foreach \i [evaluate=\i as \j using int(\i+1)] in {1,...,4} \draw[-{Latex[length=3mm]},fmmuted,line width=1.5pt] (s\i) -- (s\j);
\end{tikzpicture}
\end{center}}

\newcommand{\lo}{\mathopen{]}}\newcommand{\ro}{\mathclose{[}}
% punts ple i buit, nombres dels eixos i taula de valors
\newcommand{\pple}[2][fmblue]{\fill[#1] #2 circle (3.6pt);}
\newcommand{\nums}[4]{\foreach \k in {#1,...,#2} {\ifnum\k=0\else\node[below,font=\scriptsize,text=fmmuted] at (\k,0) {$\k$};\fi}
  \foreach \k in {#3,...,#4} {\ifnum\k=0\else\node[left,font=\scriptsize,text=fmmuted] at (0,\k) {$\k$};\fi}}
\newcommand{\tv}[3][fmblue]{\taulaclip{#1,line width=1.2pt}{\renewcommand{\arraystretch}{1.45}%
  \begin{tabular}{>{\columncolor{#1}\color{white}\heavy}c*{12}{w{c}{0.8cm}}}#2\\\rowcolor{#1!8}\cellcolor{#1}#3\end{tabular}}}
% etiqueta d'una recta o corba
\newcommand{\etq}[3]{\node[fill=white,inner sep=1.5pt,text=#1,font=\small] at #2 {#3};}


\def\intvdot#1#2#3{\ifx#2o\filldraw[fill=white,draw=#3,line width=1.3pt] (#1,0) circle (3.3pt);\fi\ifx#2c\fill[#3] (#1,0) circle (3.5pt);\fi}
\newcommand{\intv}[5]{\begin{tikzpicture}[x=0.4cm,baseline=-3pt]
  \draw[fmmuted,line width=0.8pt] (-6,0)--(7,0);
  \foreach \k in {-5,...,6} \draw[fmmuted,line width=0.6pt] (\k,-0.08)--(\k,0.08);
  \ifx#3a\def\xa{-6}\else\def\xa{#1}\fi
  \ifx#4a\def\xb{7}\else\def\xb{#2}\fi
  \ifx#4a\draw[#5,line width=2.6pt,-{Latex[length=2.6mm]}] (\xa,0)--(\xb,0);\else\ifx#3a\draw[#5,line width=2.6pt,{Latex[length=2.6mm]}-] (\xa,0)--(\xb,0);\else\draw[#5,line width=2.6pt] (\xa,0)--(\xb,0);\fi\fi
  \intvdot{#1}{#3}{#5}\intvdot{#2}{#4}{#5}
  \ifx#3a\else\node[below=2pt,font=\scriptsize] at (#1,0) {$#1$};\fi
  \ifx#4a\else\node[below=2pt,font=\scriptsize] at (#2,0) {$#2$};\fi
\end{tikzpicture}}
\newcommand{\pbu}[2][fmblue]{\draw[#1,fill=white,line width=1.2pt] #2 circle (3.4pt);}
\newcommand{\miniax}[4]{\draw[-{Latex[length=1.5mm]},fmmuted,line width=0.6pt] (#1,0)--(#2,0); \draw[-{Latex[length=1.5mm]},fmmuted,line width=0.6pt] (0,#3)--(0,#4);}
\tikzset{asim/.style={fmR!75,dashed,line width=0.9pt}}
\newcommand{\cas}[1]{\left\{\begin{array}{@{}l@{\quad}l@{}}#1\end{array}\right.}




\begin{document}
\portada{Rectes, paràboles, arrels, hipèrboles i exponencials: la fórmula i la gràfica.}

{\color{fmnavy}\heavy\Large\faMapSigns\ \ Què aprendràs?}\par\smallskip
\begin{minipage}[t]{0.485\linewidth}
\itile{sec1}{7.1}{\faChartLine}{Funcions lineals}
\itile{sec2}{7.2}{\faRoute}{Equació punt-pendent}
\itile{sec3}{7.3}{\faIcon{bezier-curve}}{Funció quadràtica}
\itile{sec4}{7.4}{\faPencilRuler}{Representació de paràboles}
\itile{sec5}{7.5}{\faIcon{square-root-alt}}{Funcions radicals}
\end{minipage}\hfill
\begin{minipage}[t]{0.485\linewidth}
\itile{sec6}{7.6}{\faBalanceScale}{Proporcionalitat inversa}
\itile{sec7}{7.7}{\faIcon{arrows-alt}}{Translacions de la hipèrbola}
\itile{sec8}{7.8}{\faRocket}{Funcions exponencials}
\itile{sec9}{7.9}{\faLightbulb}{Problemes amb funcions}
\itile{resum}{\phantom{7.9}}{\faClipboardList}{En resum}
\end{minipage}
\llegenda

% ============================================================ 7.1
\clearpage
\seccio{sec1}{Funcions lineals}{La gràfica és una recta.}

\vfill
\begin{idea}[Funció lineal \ $f(x)=mx+n$]
\textbf{Pendent} $m$: el que puja o baixa quan $x$ augmenta 1. \ \textbf{Ordenada en l'origen} $n$: talla $OY$ en $(0,n)$.\\[3pt]
$m>0$ creixent, $m<0$ decreixent. Com més gran és $|m|$, més inclinada. Mateix pendent: \textbf{paral·leles}.\\[3pt]
És contínua, sense màxims ni mínims. Talla $OX$ en $x=-\frac nm$ (si $m\neq0$).
\end{idea}

\vfill
\begin{tcbraster}[raster columns=3,raster equal height,raster column skip=8pt]
\begin{tcolorbox}[fmbase,colback=fmdefbg,colframe=fmdef,coltitle=white,title={Cas general},before upper=\centering]
\begin{tikzpicture}[scale=0.4]\miniax{-2.5}{3}{-1.5}{3.5}\draw[fmblue,line width=1.4pt] (-2.3,-1.3)--(2.5,3.5);\pple[fmR]{(0,1)}\end{tikzpicture}\par
$y=mx+n$\\ \small passa per $(0,n)$; \ $D=R=\mathbb R$
\end{tcolorbox}
\begin{tcolorbox}[fmbase,colback=fmdefbg,colframe=fmdef,coltitle=white,title={Proporcionalitat directa},before upper=\centering]
\begin{tikzpicture}[scale=0.4]\miniax{-2.5}{3}{-1.5}{3.5}\draw[fmblue,line width=1.4pt] (-1,-1.5)--(2.33,3.5);\pple[fmR]{(0,0)}\end{tikzpicture}\par
$y=mx$ \ ($n=0$)\\ \small passa per l'origen
\end{tcolorbox}
\begin{tcolorbox}[fmbase,colback=fmdefbg,colframe=fmdef,coltitle=white,title={Constant},before upper=\centering]
\begin{tikzpicture}[scale=0.4]\miniax{-2.5}{3}{-1.5}{3.5}\draw[fmblue,line width=1.4pt] (-2.3,2)--(2.6,2);\pple[fmR]{(0,2)}\end{tikzpicture}\par
$y=n$ \ ($m=0$)\\ \small horitzontal; \ $R=\{n\}$
\end{tcolorbox}
\end{tcbraster}

\vfill
\begin{exemple}[Exemple: \ representa $f(x)=-2x+3$]
\begin{minipage}[c]{0.36\linewidth}\centering
\begin{tikzpicture}[scale=0.4]
  \eixos{-2}{4}{-4}{4}
  \draw[fmblue,line width=1.5pt] (-0.4,3.8)--(3.4,-3.8);
  \draw[fmR,line width=1pt] (0,3)--(1,3)--(1,1); \node[above,font=\scriptsize,text=fmR] at (0.5,3) {1}; \node[right,font=\scriptsize,text=fmR] at (1,2) {$-2$};
  \pple[fmnavy]{(0,3)} \pple[fmnavy]{(1.5,0)}
\end{tikzpicture}
\end{minipage}\hfill
\begin{minipage}[c]{0.62\linewidth}
Dos punts n'hi ha prou: \ {\small\setlength{\tabcolsep}{4pt}\begin{tabular}{c|ccc}$x$&$0$&$1$&$3$\\\hline $y$&$3$&$1$&$-3$\end{tabular}}\\[4pt]
O bé: eixim de $(0,3)$ i, com que $m=-2$, per cada pas a la dreta baixem 2.\\[4pt]
Decreixent. \ $OY$: $(0,3)$; \ $OX$: $-2x+3=0\ \to\ \left(\frac32,0\right)$
\end{minipage}
\end{exemple}

\vfill
\begin{tcolorbox}[enhanced,colback=fmyellow!18,colframe=fmgold,arc=4mm,boxrule=1.6pt,before upper=\centering,top=2mm,bottom=2mm]
Per dos punts: \quad {\large$m=\dfrac{y_2-y_1}{x_2-x_1}$} \qquad\qquad {\large$y=y_1+m\,(x-x_1)$}
\end{tcolorbox}

\vfill
\begin{exemple}[Exemples]
Per $(-1,5)$ i $(2,-1)$: \ $m=\frac{-1-5}{2+1}=-2$, \ $y=5-2(x+1)=\boxed{-2x+3}$\\[4pt]
Pendent $\frac12$ per $(4,3)$: \ $y=3+\frac12(x-4)=\boxed{\frac12x+1}$ \qquad $OY$: $(0,1)$, \ $OX$: $(-2,0)$
\end{exemple}

\vfill

\vfill
% ============================================================ 7.2
\clearpage
\seccio{sec2}{Equació punt-pendent}{Dos punts, o un punt i el pendent, determinen una recta.}

\vfill
\begin{minipage}[c]{0.52\linewidth}
\begin{idea}[Pendent a partir de dos punts]
Si passa per $A(x_1,\,y_1)$ i $B(x_2,\,y_2)$:\\[3pt]
{\centering\Large$m=\dfrac{\cR{y_2-y_1}}{\cK{x_2-x_1}}$\par}\smallskip
el que \cR{puja la $y$} entre el que \cK{avança la $x$}.
\end{idea}
\end{minipage}\hfill
\begin{minipage}[c]{0.45\linewidth}\centering
\begin{tikzpicture}[scale=0.68]
  \eixos{-1}{5}{-1}{4}
  \draw[fmblue,line width=1.6pt] (-0.75,0.5) -- (4.5,4);
  \draw[fmK,line width=1.2pt] (0,1) -- (3,1); \draw[fmR,line width=1.2pt] (3,1) -- (3,3);
  \pple[fmnavy]{(0,1)} \pple[fmnavy]{(3,3)}
  \node[above left,font=\small\heavy] at (0,1) {$A$}; \node[above left,font=\small\heavy] at (3,3) {$B$};
  \node[below,font=\small,text=fmK!80!black] at (1.5,1) {3}; \node[right,font=\small,text=fmR] at (3,2) {2};
  \node[font=\small,text=fmblue,fill=white,inner sep=1pt] at (4.6,2.6) {$m=\frac23$};
\end{tikzpicture}
\end{minipage}

\vfill
\begin{tcolorbox}[enhanced,colback=fmyellow!18,colframe=fmgold,arc=4mm,boxrule=1.6pt,before upper=\centering,top=2mm,bottom=2mm]
{\Large Equació punt-pendent: \quad $y-y_0=m\,(x-x_0)$}\par\smallskip
la recta de pendent $m$ que passa per $(x_0,\,y_0)$
\end{tcolorbox}

\vfill
\begin{minipage}[c]{0.4\linewidth}\raggedright
\textbf{Recta per dos punts}:\\[4pt]
\pas{1} Calculem el \textbf{pendent}.\\[4pt]
\pas{2} Punt-pendent amb \textbf{un} dels punts.\\[4pt]
\pas{3} \textbf{Aïllem} la $y$.\\[4pt]
\pas{4} \textbf{Comprovem} amb l'altre punt.
\end{minipage}\hfill
\begin{minipage}[c]{0.57\linewidth}
\begin{exemple}[Exemple: \ recta per $A(1,\,-2)$ i $B(4,\,7)$]
\pas{1} $m=\dfrac{7-(-2)}{4-1}=\dfrac93=3$\\[4pt]
\pas{2} $y-(-2)=3\,(x-1)$\\[4pt]
\pas{3} $y+2=3x-3\ \to\ \boxed{y=3x-5}$\\[4pt]
\pas{4} \nota{$B$: \ $3\cdot4-5=7$ \be{}}
\end{exemple}
\end{minipage}

\vfill
\begin{tcbraster}[raster columns=2,raster equal height,raster column skip=10pt]
\begin{tcolorbox}[fmbase,colback=fmdefbg,colframe=fmdef,coltitle=white,title={Forma general: \ $ax+by+c=0$}]
Per a passar a $y=mx+n$, \textbf{aïllem la $y$}:\par\smallskip
{\centering$4x-2y+6=0\ \to\ \boxed{y=2x+3}$\par}
\end{tcolorbox}
\begin{tcolorbox}[fmbase,colback=fmdefbg,colframe=fmdef,coltitle=white,title={Rectes paral·leles: mateix pendent}]
Paral·lela a $y=-2x+1$ per $(3,\,4)$:\par\smallskip
{\centering$y-4=-2\,(x-3)\ \to\ \boxed{y=-2x+10}$\par}
\end{tcolorbox}
\end{tcbraster}

\vfill
\begin{exemple}[Exemple: \ recta per $(-2,\,7)$ i $(4,\,-2)$, i tall amb l'eix $Y$]
\begin{minipage}[t]{0.5\linewidth}
$m=\dfrac{-2-7}{4-(-2)}=\dfrac{-9}{6}=-\dfrac32$\\[4pt]
$y-7=-\tfrac32\,(x+2)\ \to\ \boxed{y=-\tfrac32x+4}$
\end{minipage}\hfill
\begin{minipage}[t]{0.47\linewidth}
Talla l'eix $Y$ en $(0,\,4)$.\\[4pt]
\nota{comprovació: $-\tfrac32\cdot4+4=-2$ \be{}}
\end{minipage}
\end{exemple}

\vfill
\begin{exemple}[Exemple: \ estan alineats $A(-1,\,-4)$, $B(1,\,2)$ i $C(3,\,7)$?]
$m_{AB}=\dfrac{2-(-4)}{1-(-1)}=\dfrac62=3$ \qquad $m_{BC}=\dfrac{7-2}{3-1}=\dfrac52$ \qquad $3\neq\dfrac52$: \ \textbf{no} estan alineats.
\end{exemple}

\vfill
% ============================================================ 7.3
\clearpage
\seccio{sec3}{Funció quadràtica: la paràbola $y=ax^2$}{Un polinomi de segon grau dibuixa una paràbola.}

\vfill
\begin{idea}[Funció quadràtica]
{\centering\Large$y=ax^2+bx+c\qquad (a\neq0)$\par}\smallskip
La seua gràfica és una corba simètrica anomenada \textbf{paràbola}. La més senzilla és $y=x^2$.
\end{idea}

\vfill
\begin{minipage}[c]{0.5\linewidth}
{\centering\tv{$x$ & $-3$ & $-2$ & $-1$ & 0 & 1 & 2 & 3}{$y$ & 9 & 4 & 1 & 0 & 1 & 4 & 9}\par}\medskip
\faCheck\ \textbf{Simètrica} respecte de l'eix $Y$: $x$ i $-x$ tenen la mateixa imatge. Aquesta recta és l'\textbf{eix de simetria}.\\[4pt]
\faCheck\ El \textbf{vèrtex} és $V(0,\,0)$: el \textbf{mínim}.\\[4pt]
\faCheck\ Decreix en $\lo-\infty,\,0\ro$ i creix en $\lo0,\,+\infty\ro$.
\end{minipage}\hfill
\begin{minipage}[c]{0.46\linewidth}\centering
\begin{tikzpicture}[x=0.6cm,y=0.36cm]
  \eixos{-3}{3}{-1}{9}
  \nums{-3}{3}{1}{9}
  \draw[fmK,line width=1pt,dashed] (0,-1) -- (0,9.3);
  \draw[fmblue,line width=1.7pt,domain=-3.05:3.05,samples=50] plot (\x,{\x*\x});
  \foreach \x in {-3,...,3} \fill[fmR] (\x,{\x*\x}) circle (3.5pt);
  \node[below right,font=\small\heavy,text=fmR] at (0,0) {$V$};
\end{tikzpicture}
\end{minipage}

\vfill
\subsection{Què fa el coeficient $a$?}
\begin{minipage}[c]{0.52\linewidth}
\chip{fmK}{$a>0$} \ branques cap \textbf{amunt} $\cup$: el vèrtex és un \textbf{mínim}.\\[5pt]
\chip{fmR}{$a<0$} \ branques cap \textbf{avall} $\cap$: el vèrtex és un \textbf{màxim}.\\[5pt]
\chip{fmblue}{$|a|$ gran} \ més \textbf{tancada}; \ \chip{fmblue}{$|a|<1$} \ més \textbf{oberta} que $y=x^2$.
\end{minipage}\hfill
\begin{minipage}[c]{0.22\linewidth}\centering
\begin{tikzpicture}[scale=0.32]
  \eixos{-3}{3}{-1}{8}
  \draw[fmK,line width=1.5pt,domain=-1.6:1.6,samples=40] plot (\x,{3*\x*\x});
  \draw[fmK!55,line width=1.5pt,domain=-3:3,samples=40] plot (\x,{\x*\x/3});
\end{tikzpicture}\\{\scriptsize\color{fmK!80!black}$3x^2$ \ i \ $\frac13x^2$}
\end{minipage}\hfill
\begin{minipage}[c]{0.22\linewidth}\centering
\begin{tikzpicture}[scale=0.32]
  \eixos{-3}{3}{-8}{1}
  \draw[fmR,line width=1.5pt,domain=-2.8:2.8,samples=40] plot (\x,{-\x*\x});
  \draw[fmR!55,line width=1.5pt,domain=-3:3,samples=40] plot (\x,{-\x*\x/4});
\end{tikzpicture}\\{\scriptsize\color{fmR}$-x^2$ \ i \ $-\frac14x^2$}
\end{minipage}

\vfill
\begin{minipage}[c]{0.6\linewidth}
\begin{tcolorbox}[fmbase,colback=fmdefbg,colframe=fmdef,coltitle=white,title={Desplaçament: \ $y=ax^2+c$}]
És la paràbola $y=ax^2$ desplaçada \textbf{$c$ unitats} cap amunt (o avall si $c<0$). El vèrtex és $(0,\,c)$.\par\smallskip
{\centering$y=x^2-4$: \ vèrtex $(0,\,-4)$, talls $x^2=4\to x=\pm2$\par}
\end{tcolorbox}
\end{minipage}\hfill
\begin{minipage}[c]{0.36\linewidth}\centering
\begin{tikzpicture}[scale=0.38]
  \eixos{-3}{3}{-4}{5}
  \draw[fmmuted,line width=1.2pt,dashed,domain=-2.2:2.2,samples=40] plot (\x,{\x*\x});
  \draw[fmblue,line width=1.6pt,domain=-3:3,samples=40] plot (\x,{\x*\x-4});
  \fill[fmR] (0,-4) circle (5pt); \fill[fmK] (-2,0) circle (5pt); \fill[fmK] (2,0) circle (5pt);
\end{tikzpicture}
\end{minipage}

\vfill
\begin{exemple}[Exemple: \ $a$ perquè $y=ax^2$ passe per $(2,\,-6)$]
$-6=a\cdot2^2\ \to\ \boxed{a=-\tfrac32}$ \qquad Passa també per $(-2,\,-6)$, per \textbf{simetria}: $x=-2$ té la mateixa imatge que $x=2$.
\end{exemple}

\vfill
% ============================================================ 7.4
\clearpage
\seccio{sec4}{Representació de paràboles}{La gràfica és una paràbola.}

\vfill
\begin{idea}[Funció quadràtica \ $f(x)=ax^2+bx+c$, \ $a\neq0$]
$a>0$: branques cap amunt $\cup$, el vèrtex és el \textbf{mínim}. \ $a<0$: cap avall $\cap$, el vèrtex és el \textbf{màxim}. Com més gran $|a|$, més tancada.\\[3pt]
\textbf{Vèrtex}: \ $V_x=\dfrac{-b}{2a}$, \ $V_y=f(V_x)$. \ \textbf{Eix de simetria}: $x=V_x$.\\[3pt]
Talls: $OY$ en $(0,c)$; $OX$, les solucions de $ax^2+bx+c=0$: dos, un o cap, segons $\Delta=b^2-4ac$.
\end{idea}

\vfill
\begin{exemple}[Exemple: \ representa $f(x)=x^2+2x-3$]
\begin{minipage}[c]{0.4\linewidth}\centering
\begin{tikzpicture}[scale=0.4]
  \eixos{-5}{3}{-5}{5}
  \draw[fmmuted,dashed] (-1,-5)--(-1,5);
  \draw[fmblue,line width=1.5pt,domain=-4.4:2.4,smooth,samples=40] plot (\x,{\x*\x+2*\x-3});
  \pple[fmR]{(-1,-4)} \pple[fmnavy]{(1,0)} \pple[fmnavy]{(-3,0)} \pple[fmnavy]{(0,-3)}
  \node[below,font=\scriptsize\heavy,text=fmR] at (-1,-4.1) {$V$};
\end{tikzpicture}
\end{minipage}\hfill
\begin{minipage}[c]{0.58\linewidth}
\pas{1} Forma: \ $a=1>0$ \ $\to$ \ $\cup$\\[4pt]
\pas{2} Talls: \ $OY$ $(0,-3)$; \ $x^2+2x-3=0\ \to\ (1,0),\ (-3,0)$\\[4pt]
\pas{3} Vèrtex: \ $V_x=\frac{-2}{2}=-1$, \ $V_y=f(-1)=-4$ \ $\to$ \ $V(-1,-4)$\\[4pt]
\pas{4} Unim els punts amb una paràbola simètrica respecte de $x=-1$. Si cal, taula de valors.
\end{minipage}
\end{exemple}

\vfill
\begin{tcbraster}[raster columns=3,raster equal height,raster column skip=8pt]
\begin{tcolorbox}[fmbase,colback=fmB!8,colframe=fmB,coltitle=white,title={$y=ax^2$}]
Vèrtex $(0,0)$, eix $x=0$.\\ \small$R=[0,+\infty\ro$ si $a>0$
\end{tcolorbox}
\begin{tcolorbox}[fmbase,colback=fmB!8,colframe=fmB,coltitle=white,title={$y=ax^2+c$}]
Vèrtex $(0,c)$, eix $x=0$.\\ \small$R=[c,+\infty\ro$ si $a>0$
\end{tcolorbox}
\begin{tcolorbox}[fmbase,colback=fmB!8,colframe=fmB,coltitle=white,title={$y=ax^2+bx$}]
Passa per $(0,0)$.\\ L'altre tall: $x=-\frac ba$.
\end{tcolorbox}
\end{tcbraster}

\vfill
\begin{truc}
Si ve \textbf{factoritzada}, $y=a(x-x_1)(x-x_2)$, els talls amb $OX$ són $x_1$ i $x_2$, i el vèrtex és al mig: $V_x=\dfrac{x_1+x_2}{2}$.
\end{truc}

\vfill
\begin{exemple}[Exemples]
$y=-x^2+4x-3$: \ $\cap$; \ $OX$ $(1,0)$, $(3,0)$; \ $OY$ $(0,-3)$; \ $V(2,1)$; \ $R=\lo-\infty,1]$\\[4pt]
$y=2x^2-8$: \ $\cup$; \ $OX$ $(-2,0)$, $(2,0)$; \ $V=(0,-8)$; \ $R=[-8,+\infty\ro$\\[4pt]
$y=x^2-6x+10$: \ $\cup$; \ $\Delta=-4<0$, no talla $OX$; \ $OY$ $(0,10)$; \ $V(3,1)$; \ taula: $f(1)=f(5)=5$
\end{exemple}

\vfill

\vfill
% ============================================================ 7.5
\clearpage
\seccio{sec5}{Funcions radicals}{La $x$ dins d'una arrel: només hi entren els nombres que no la fan negativa.}

\vfill
\begin{minipage}[c]{0.42\linewidth}\centering
\begin{tikzpicture}[scale=0.55]
  \eixos{-1}{9}{-1}{4}
  \draw[fmblue,line width=1.8pt,domain=0:9,samples=60] plot (\x,{sqrt(\x)});
  \foreach \p in {(0,0),(1,1),(4,2),(9,3)} \fill[fmblue] \p circle (4pt);
  \node[fmblue,font=\small\heavy] at (6,3.3) {$y=\sqrt x$};
\end{tikzpicture}\\[3pt]
\tv{$x$&0&1&4&9}{$\sqrt x$&0&1&2&3}
\end{minipage}\hfill
\begin{minipage}[c]{0.55\linewidth}
\begin{idea}[La funció \ $f(x)=\sqrt x$]
L'arrel d'un negatiu no existeix: la $x$ només pot ser positiva o zero.\\[3pt]
$D=[0,+\infty\ro$, \ $R=[0,+\infty\ro$.\\[3pt]
Comença en $(0,0)$ i va cap a la dreta i cap amunt: \textbf{creixent}, cada vegada més a poc a poc. És \textbf{mitja paràbola} tombada.
\end{idea}
\begin{truc}[Taula fàcil]
Dona a la $x$ valors que facen el radicand un \textbf{quadrat perfecte}: 0, 1, 4, 9\dots
\end{truc}
\end{minipage}

\vfill
\begin{minipage}[c]{0.5\linewidth}\centering\small
\taularod{>{\heavy}l<{\hspace{3pt}}>{\hspace{3pt}}c<{\hspace{3pt}}>{\hspace{3pt}}l<{\hspace{3pt}}}{
\rowcolor{fmblue}\color{white}Translació & \color{white}\heavy Funció & \color{white}\heavy Comença en\\
Vertical & $\sqrt x\pm b$ & $(0,\,\pm b)$\\
\rowcolor{fmsurface}Horitzontal & $\sqrt{x\pm a}$ & $(\mp a,\,0)$ \ \nota{al revés!}\\
Les dues & $\sqrt{x\pm a}\pm b$ & $(\mp a,\,\pm b)$\\
\rowcolor{fmsurface}Signe menys & $-\sqrt x$ & $(0,0)$, cap \textbf{avall}
}
\end{minipage}\hfill
\begin{minipage}[c]{0.47\linewidth}\centering
\begin{tikzpicture}[scale=0.42]
  \eixos{-4}{6}{-3}{5}
  \draw[fmmuted,dashed,line width=1pt,domain=0:6,samples=40] plot (\x,{sqrt(\x)});
  \draw[fmA,line width=1.5pt,domain=0:6,samples=40] plot (\x,{sqrt(\x)+3}) node[right,font=\scriptsize\heavy] {$\sqrt x+3$};
  \draw[fmB,line width=1.5pt,domain=-3:6,samples=60] plot (\x,{sqrt(\x+3)}) node[right,font=\scriptsize\heavy] {$\sqrt{x+3}$};
  \draw[fmR,line width=1.5pt,domain=0:6,samples=40] plot (\x,{-sqrt(\x)}) node[right,font=\scriptsize\heavy] {$-\sqrt x$};
\end{tikzpicture}
\end{minipage}

\vfill
\begin{exemple}[Exemple: \ representa $f(x)=\sqrt{x+2}-3$]
\begin{minipage}[c]{0.36\linewidth}\centering
\begin{tikzpicture}[scale=0.38]
  \eixos{-3}{8}{-4}{2}
  \draw[fmblue,line width=1.6pt,domain=-2:8,samples=60] plot (\x,{sqrt(\x+2)-3});
  \foreach \p in {(-2,-3),(-1,-2),(2,-1),(7,0)} \fill[fmblue] \p circle (5pt);
\end{tikzpicture}
\end{minipage}\hfill
\begin{minipage}[c]{0.62\linewidth}
\pas{1} Domini: \ $x+2\ge0\ \to\ D=[-2,+\infty\ro$\\[3pt]
\pas{2} Punt inicial: \ $f(-2)=-3$: \ $(-2,-3)$ \nota{2 a l'esquerra i 3 avall}\\[3pt]
\pas{3} {\small\setlength{\tabcolsep}{3pt}\begin{tabular}{c|cccc}$x$&$-2$&$-1$&$2$&$7$\\\hline $y$&$-3$&$-2$&$-1$&$0$\end{tabular}}\\[3pt]
$R=[-3,+\infty\ro$, creixent. Talls: $OX$ $(7,0)$, \ $OY$ $(0,\sqrt2-3)\approx(0;-1{,}59)$
\end{minipage}
\end{exemple}

\vfill
\begin{exemple}[Exemples: \ on comença, domini i recorregut]
$y=\sqrt{x+1}$: \ $(-1,0)$, \ $D=[-1,+\infty\ro$, \ $R=[0,+\infty\ro$ \qquad $y=\sqrt x-4$: \ $(0,-4)$, \ $R=[-4,+\infty\ro$\\[4pt]
$y=\sqrt{x-2}+1$: \ $(2,1)$, \ $D=[2,+\infty\ro$, \ $R=[1,+\infty\ro$ \qquad $y=-\sqrt{x+1}$: \ $(-1,0)$, \ $R=\lo-\infty,0]$ \nota{cap avall}
\end{exemple}

\vfill
% ============================================================ 7.6
\clearpage
\seccio{sec6}{Funcions de proporcionalitat inversa}{El producte és constant: $x\cdot y=k$.}

\vfill
\begin{minipage}[c]{0.44\linewidth}
{\small Recórrer 120 km: \ $v\cdot t=120$, \ $t=\dfrac{120}{v}$}\par\smallskip
\tv{$v$ (km/h)&$20$&$40$&$60$&$120$}{$t$ (h)&$6$&$3$&$2$&$1$}\\[3pt]
{\comic\small\color{fmmuted}el doble de velocitat, la meitat de temps}
\end{minipage}\hfill
\begin{minipage}[c]{0.53\linewidth}
\begin{idea}[Funció \ $f(x)=\dfrac kx$, \ $k\neq0$]
Gràfica: \textbf{hipèrbola} de dues branques. No està definida en $x=0$ i mai no val 0.\\[2pt]
\textbf{Asímptotes} (rectes a què s'acosta sense tocar-les): els eixos, $x=0$ i $y=0$.\\[2pt]
$D=R=\mathbb R-\{0\}$. Sense màxims ni mínims; no talla els eixos.
\end{idea}
\end{minipage}

\vfill
\begin{tcbraster}[raster columns=2,raster equal height,raster column skip=10pt]
\begin{tcolorbox}[fmbase,colback=fmR!6,colframe=fmR,coltitle=white,title={$k>0$: \ decreixent a cada branca}]
\begin{minipage}[c]{0.45\linewidth}\centering
\begin{tikzpicture}[scale=0.38]\miniax{-4}{4}{-4}{4}
\draw[fmblue,line width=1.4pt,domain=0.8:3.8,smooth] plot (\x,{3/\x}); \draw[fmblue,line width=1.4pt,domain=-3.8:-0.8,smooth] plot (\x,{3/\x});\end{tikzpicture}
\end{minipage}\hfill
\begin{minipage}[c]{0.52\linewidth}$y=\dfrac3x$\\[3pt] \small 1r i 3r quadrant\end{minipage}
\end{tcolorbox}
\begin{tcolorbox}[fmbase,colback=fmK!8,colframe=fmK,coltitle=white,title={$k<0$: \ creixent a cada branca}]
\begin{minipage}[c]{0.45\linewidth}\centering
\begin{tikzpicture}[scale=0.38]\miniax{-4}{4}{-4}{4}
\draw[fmblue,line width=1.4pt,domain=0.3:3.8,smooth] plot (\x,{-1/\x}); \draw[fmblue,line width=1.4pt,domain=-3.8:-0.3,smooth] plot (\x,{-1/\x});\end{tikzpicture}
\end{minipage}\hfill
\begin{minipage}[c]{0.52\linewidth}$y=-\dfrac1x$\\[3pt] \small 2n i 4t quadrant\end{minipage}
\end{tcolorbox}
\end{tcbraster}

\vfill
\begin{exemple}[Exemple: \ representa $f(x)=\dfrac4x$]
\begin{minipage}[c]{0.36\linewidth}\centering
\begin{tikzpicture}[x=0.3cm,y=0.17cm]
  \eixos{-5}{5}{-8}{8}
  \draw[fmblue,line width=1.5pt,domain=0.5:5,smooth,samples=40] plot (\x,{4/\x});
  \draw[fmblue,line width=1.5pt,domain=-5:-0.5,smooth,samples=40] plot (\x,{4/\x});
  \foreach \p in {(1,4),(2,2),(4,1),(-1,-4),(-2,-2),(-4,-1)} \fill[fmblue] \p circle (5pt);
\end{tikzpicture}
\end{minipage}\hfill
\begin{minipage}[c]{0.62\linewidth}
Valors a \textbf{banda i banda} de $x=0$, alguns a prop ($\pm\frac12$):\\[3pt]
{\small\setlength{\tabcolsep}{3pt}\begin{tabular}{c|cccccccc}$x$&$-4$&$-2$&$-1$&$-\frac12$&$\frac12$&$1$&$2$&$4$\\\hline $y$&$-1$&$-2$&$-4$&$-8$&$8$&$4$&$2$&$1$\end{tabular}}\\[4pt]
$k=4>0$: 1r i 3r quadrant, decreixent a cada branca.
\end{minipage}
\end{exemple}

\vfill
\begin{compte}[«A cada branca», no «sempre»]
$f(-1)=-4<f(1)=4$ encara que $-1<1$: la funció no és decreixent en tot el domini, sinó \textbf{a cada branca}.
\end{compte}

\vfill
\begin{exemple}[Exemples]
La hipèrbola $y=\frac kx$ passa per $(3,-4)$: \ $k=x\cdot y=-12$, \ $y=-\dfrac{12}{x}$ \qquad Com més gran és $|k|$, més lluny de l'origen.\\[4pt]
Un premi a parts iguals: amb 4 persones, 150 € cadascuna. \ $k=600$, \ $y=\dfrac{600}{x}$; \ amb 6 persones, \boxed{100\text{ €}}
\end{exemple}

\vfill
% ============================================================ 7.7
\clearpage
\seccio{sec7}{Translacions de la hipèrbola}{La forma no canvia: les asímptotes es mouen amb la gràfica.}

\vfill
\begin{minipage}[c]{0.55\linewidth}\centering\small
\taularod{>{\heavy}l<{\hspace{3pt}}>{\hspace{3pt}}c<{\hspace{3pt}}>{\hspace{3pt}}l<{\hspace{3pt}}}{
\rowcolor{fmblue}\color{white}Translació & \color{white}\heavy Funció & \color{white}\heavy Asímptotes\\
Vertical & $\dfrac kx\pm b$\rule[-9pt]{0pt}{24pt} & $x=0$, \ $y=\pm b$\\
\rowcolor{fmsurface}Horitzontal & $\dfrac k{x\pm a}$\rule[-9pt]{0pt}{24pt} & $x=\mp a$ \nota{al revés!}, \ $y=0$\\
Les dues & $\dfrac k{x\pm a}\pm b$\rule[-9pt]{0pt}{24pt} & $x=\mp a$, \ $y=\pm b$
}
\end{minipage}\hfill
\begin{minipage}[c]{0.42\linewidth}
\begin{tcolorbox}[enhanced,colback=fmyellow!18,colframe=fmgold,arc=4mm,boxrule=1.6pt,before upper=\centering,top=2mm,bottom=2mm]
{\Large$y=\dfrac{k}{x-a}+b$}\par\smallskip
\small vertical $x=a$ (on s'anul·la el denominador), \ horitzontal $y=b$ (el que sumem)\\[2pt]
$D=\mathbb R-\{a\}$, \ $R=\mathbb R-\{b\}$
\end{tcolorbox}
\end{minipage}

\vfill
\begin{tcbraster}[raster columns=3,raster equal height,raster column skip=8pt]
\begin{tcolorbox}[fmbase,colback=fmdefbg,colframe=fmdef,coltitle=white,title={$y=\frac3x+2$},before upper=\centering]
\begin{tikzpicture}[scale=0.28]\miniax{-5}{5}{-4}{7}\draw[asim] (-5,2)--(5,2);
\draw[fmblue,line width=1.3pt,domain=0.5:5,smooth] plot (\x,{3/\x+2}); \draw[fmblue,line width=1.3pt,domain=-5:-0.6,smooth] plot (\x,{3/\x+2});\end{tikzpicture}\par
\small puja 2
\end{tcolorbox}
\begin{tcolorbox}[fmbase,colback=fmdefbg,colframe=fmdef,coltitle=white,title={$y=\frac3{x-1}$},before upper=\centering]
\begin{tikzpicture}[scale=0.28]\miniax{-5}{5}{-5}{5}\draw[asim] (1,-5)--(1,5);
\draw[fmblue,line width=1.3pt,domain=1.6:5,smooth] plot (\x,{3/(\x-1)}); \draw[fmblue,line width=1.3pt,domain=-5:0.4,smooth] plot (\x,{3/(\x-1)});\end{tikzpicture}\par
\small 1 a la dreta
\end{tcolorbox}
\begin{tcolorbox}[fmbase,colback=fmdefbg,colframe=fmdef,coltitle=white,title={$y=\frac3{x+3}-2$},before upper=\centering]
\begin{tikzpicture}[scale=0.28]\miniax{-7}{3}{-6}{3}\draw[asim] (-3,-6)--(-3,3); \draw[asim] (-7,-2)--(3,-2);
\draw[fmblue,line width=1.3pt,domain=-2.4:3,smooth] plot (\x,{3/(\x+3)-2}); \draw[fmblue,line width=1.3pt,domain=-7:-3.7,smooth] plot (\x,{3/(\x+3)-2});\end{tikzpicture}\par
\small 3 a l'esquerra i 2 avall
\end{tcolorbox}
\end{tcbraster}

\vfill
\begin{exemple}[Exemple: \ representa $f(x)=\dfrac{3}{x-2}+1$]
\begin{minipage}[c]{0.38\linewidth}\centering
\begin{tikzpicture}[scale=0.3]
  \eixos{-3}{7}{-6}{8}
  \draw[asim] (2,-6)--(2,8); \draw[asim] (-3,1)--(7,1);
  \draw[fmblue,line width=1.5pt,domain=-3:1.57,smooth,samples=40] plot (\x,{3/(\x-2)+1});
  \draw[fmblue,line width=1.5pt,domain=2.43:7,smooth,samples=40] plot (\x,{3/(\x-2)+1});
\end{tikzpicture}
\end{minipage}\hfill
\begin{minipage}[c]{0.6\linewidth}
\pas{1} Asímptotes, amb línia discontínua: \ $x=2$, \ $y=1$\\[3pt]
\pas{2} Valors a banda i banda de $x=2$:\\[2pt]
{\small\setlength{\tabcolsep}{3pt}\begin{tabular}{c|cccccc}$x$&$-1$&$1$&$1{,}5$&$2{,}5$&$3$&$5$\\\hline $y$&$0$&$-2$&$-5$&$7$&$4$&$2$\end{tabular}}\\[4pt]
$D=\mathbb R-\{2\}$, \ $R=\mathbb R-\{1\}$; \ $k>0$: decreixent a cada branca\\[3pt]
Talls: \ $OY$ $\left(0,-\frac12\right)$, \ $OX$ $(-1,0)$
\end{minipage}
\end{exemple}

\vfill
\begin{exemple}[Exemples: \ asímptotes]
$y=\frac{2}{x-3}$: \ $x=3$, $y=0$ \qquad $y=-\frac1x+2$: \ $x=0$, $y=2$, creixent \qquad $y=\frac3{x+1}-2$: \ $x=-1$, $y=-2$\\[4pt]
$y=\frac{5}{x+4}+3$: \ $x=-4$, $y=3$, \ $D=\mathbb R-\{-4\}$, \ $R=\mathbb R-\{3\}$
\end{exemple}

\vfill
% ============================================================ 7.8
\clearpage
\seccio{sec8}{Funcions exponencials}{La $x$ a l'exponent: cada pas, multiplica per $a$.}

\vfill
\begin{minipage}[c]{0.5\linewidth}
\begin{idea}[Funció \ $f(x)=a^x$, \ $a>0$, $a\neq1$]
Quan $x$ augmenta 1, la $y$ es \textbf{multiplica per $a$}.\\[3pt]
$D=\mathbb R$, \ $R=\lo0,+\infty\ro$. Passa per $(0,1)$. Asímptota horitzontal $y=0$.
\end{idea}
\end{minipage}\hfill
\begin{minipage}[c]{0.47\linewidth}\centering\small
\taularod{>{\heavy}c*{6}{>{\hspace{2pt}}c<{\hspace{2pt}}}}{%
\rowcolor{fmblue}\color{white}$x$ & \color{white}$-2$ & \color{white}$-1$ & \color{white}$0$ & \color{white}$1$ & \color{white}$2$ & \color{white}$3$\\
$2^x$ & $\frac14$ & $\frac12$ & $1$ & $2$ & $4$ & $8$\\
\rowcolor{fmsurface}$\left(\frac12\right)^x$ & $4$ & $2$ & $1$ & $\frac12$ & $\frac14$ & $\frac18$
}
\end{minipage}

\vfill
\begin{tcbraster}[raster columns=2,raster equal height,raster column skip=10pt]
\begin{tcolorbox}[fmbase,colback=fmK!8,colframe=fmK,coltitle=white,title={$a>1$: \ creixent}]
\begin{minipage}[c]{0.45\linewidth}\centering
\begin{tikzpicture}[scale=0.4]\miniax{-3.5}{3}{-0.5}{4.5}
\draw[fmblue,line width=1.4pt,domain=-3.4:2.1,smooth] plot (\x,{2^\x}); \pple[fmR]{(0,1)}\end{tikzpicture}
\end{minipage}\hfill
\begin{minipage}[c]{0.52\linewidth}$y=2^x$\\[3pt] \small asímptota $y=0$ per l'esquerra\end{minipage}
\end{tcolorbox}
\begin{tcolorbox}[fmbase,colback=fmR!6,colframe=fmR,coltitle=white,title={$0<a<1$: \ decreixent}]
\begin{minipage}[c]{0.45\linewidth}\centering
\begin{tikzpicture}[scale=0.4]\miniax{-3}{3.5}{-0.5}{4.5}
\draw[fmblue,line width=1.4pt,domain=-2.1:3.4,smooth] plot (\x,{2^(-\x)}); \pple[fmR]{(0,1)}\end{tikzpicture}
\end{minipage}\hfill
\begin{minipage}[c]{0.52\linewidth}$y=\left(\frac12\right)^x=2^{-x}$\\[3pt] \small simètrica de $2^x$ respecte de $OY$\end{minipage}
\end{tcolorbox}
\end{tcbraster}

\vfill
\begin{compte}[Mai no val 0]
$R=\lo0,+\infty\ro$, no $[0,+\infty\ro$: una potència de base positiva mai no val 0.
\end{compte}

\vfill
\begin{center}\small
\taularod{>{\heavy}l<{\hspace{4pt}}>{\hspace{4pt}}c<{\hspace{4pt}}>{\hspace{4pt}}c<{\hspace{4pt}}>{\hspace{4pt}}c<{\hspace{4pt}}}{
\rowcolor{fmblue}\color{white}Desplaçament & \color{white}\heavy Funció & \color{white}\heavy Asímptota & \color{white}\heavy Punt de referència\\
Vertical & $a^x+b$ & $y=b$ & $(0,\,1+b)$\\
\rowcolor{fmsurface}Horitzontal & $a^{x-c}$ & $y=0$ & $(c,\,1)$\\
Exemple & $3^{x-2}+1$ & $y=1$ & $(2,\,2)$
}
\end{center}

\vfill
\begin{exemple}[Exemples]
$y=2^x+1$: \ $y=1$, \ $(0,2)$, \ $R=\lo1,+\infty\ro$ \qquad $y=\left(\frac12\right)^x-3$: \ $y=-3$, \ $(0,-2)$, \ decreixent\\[4pt]
$y=2^{x+1}-2$: \ $y=-2$, \ $(-1,-1)$, \ $R=\lo-2,+\infty\ro$, \ creixent
\end{exemple}

\vfill
\begin{idea}[Creixement i decreixement exponencial: \ $y=C\cdot a^t$]
$C$: quantitat inicial. Cada unitat de temps es multiplica per $a$: \ $a>1$ creix, \ $0<a<1$ decreix (sense arribar a 0).
\end{idea}

\vfill
\begin{center}\small
\taularod{>{\raggedright\arraybackslash}p{7.2cm}<{\hspace{4pt}}>{\hspace{4pt}}c<{\hspace{4pt}}>{\hspace{4pt}}c<{\hspace{4pt}}}{
\rowcolor{fmblue}\multicolumn{1}{l}{\color{white}\heavy Situació} & \color{white}\heavy $a$ & \color{white}\heavy Funció\\
100 bacteris que es dupliquen cada hora & $2$ & $N(t)=100\cdot2^t$\\
\rowcolor{fmsurface}1\,000 € al 4\,\% d'interès compost anual & $1{,}04$ & $C(t)=1\,000\cdot1{,}04^t$\\
Cotxe de 20\,000 € que perd el 15\,\% cada any & $0{,}85$ & $V(t)=20\,000\cdot0{,}85^t$
}
\end{center}

\vfill

\vfill
% ============================================================ 7.9
\clearpage
\seccio{sec9}{Problemes amb funcions}{Quina funció hi ha darrere de la situació?}

\vfill
\begin{center}\small
\taularod{>{\raggedright\arraybackslash}p{8.6cm}<{\hspace{4pt}}>{\hspace{4pt}}l<{\hspace{4pt}}>{\hspace{4pt}}c<{\hspace{4pt}}}{
\rowcolor{fmblue}\multicolumn{1}{l}{\color{white}\heavy Pista de l'enunciat} & \color{white}\heavy Tipus & \color{white}\heavy Funció\\
Una quantitat fixa més una quantitat per unitat (per hora, per km\dots) & Lineal & $y=mx+n$\\
\rowcolor{fmsurface}Un producte (preu $\times$ unitats, base $\times$ altura), un màxim o un mínim & Quadràtica & $y=ax^2+bx+c$\\
Producte constant (més aixetes, menys temps) & P. inversa & $y=\frac kx$\\
Es multiplica pel mateix nombre cada cert temps & Exponencial & $y=C\cdot a^t$
}
\end{center}

\vfill
\begin{tcbraster}[raster columns=4,raster equal height,raster column skip=6pt]
\begin{tcolorbox}[fmbase,colback=fmdefbg,colframe=fmdef,coltitle=white,title={\pas{1} Variables}]
\small i quins valors pot prendre la $x$.
\end{tcolorbox}
\begin{tcolorbox}[fmbase,colback=fmdefbg,colframe=fmdef,coltitle=white,title={\pas{2} Funció}]
\small Escrivim la fórmula.
\end{tcolorbox}
\begin{tcolorbox}[fmbase,colback=fmdefbg,colframe=fmdef,coltitle=white,title={\pas{3} Calculem}]
\small Imatge, antiimatge o vèrtex.
\end{tcolorbox}
\begin{tcolorbox}[fmbase,colback=fmdefbg,colframe=fmdef,coltitle=white,title={\pas{4} Responem}]
\small Amb una frase i les unitats.
\end{tcolorbox}
\end{tcbraster}

\vfill
\begin{exemple}[Lineal: \ patinets]
A: 2 € de desbloqueig i 0,25 €/min. \ B: 0,45 €/min. \ 20 min? \ Quan convé A?\par\smallskip
$A(t)=0{,}25t+2$, \ $B(t)=0{,}45t$ \qquad $A(20)=7$ €, \ $B(20)=9$ €\\[3pt]
$0{,}25t+2=0{,}45t\ \to\ t=10$: \ fins a 10 min convé B; \ a partir de 10 min, \textbf{A}
\end{exemple}

\vfill
\begin{exemple}[Quadràtica: \ ingressos]
A $p$ euros se'n venen $300-5p$ unitats. \ $I(p)=p(300-5p)=-5p^2+300p$ \ ($\cap$, el vèrtex és el màxim)\\[3pt]
$p=\frac{-300}{-10}=30$, \ $I(30)=4\,500$: \ \textbf{4\,500 € a 30 €} \qquad Si $p\in[35,50]$: \ $I(35)=4\,375$, $I(50)=2\,500$: \ 35 €
\end{exemple}

\vfill
\begin{exemple}[Quadràtica: \ altura d'una pilota \ $h(t)=-5t^2+20t+1$]
Vèrtex: \ $t=\frac{-20}{-10}=2$, \ $h(2)=21$: \ \textbf{21 m als 2 s} \qquad Torna a 1 m: \ $-5t(t-4)=0\ \to\ t=4$ s
\end{exemple}

\newpage
\begin{exemple}[Proporcionalitat inversa: \ aixetes]
4 aixetes, 6 h: \ $k=24$, \ $t(n)=\frac{24}{n}$ \qquad $t(8)=3$ h \qquad $\frac{24}{n}=2\ \to\ n=12$ aixetes
\end{exemple}

\vfill
\begin{exemple}[Exponencial: \ bacteris]
500 que es tripliquen cada hora: \ $N(t)=500\cdot3^t$, \ $N(4)=40\,500$\\[3pt]
Més d'un milió: \ $3^t>2\,000$; \ $3^6=729$, \ $3^7=2\,187$: \ \textbf{7 hores}
\end{exemple}

\vfill
\begin{exemple}[Radical: \ caiguda lliure]
Una pedra tarda $t=\sqrt{\dfrac{h}{4{,}9}}$ segons a caure des de $h$ metres: \ des de 20 m, \ $t\approx2{,}02$ s; \ en 3 s cau $h=4{,}9\cdot9=44{,}1$ m
\end{exemple}

\vfill
\begin{exemple}[Més problemes]
Llanterner: 30 € + 24 €/h: \ $P(t)=24t+30$, \ $P(2{,}5)=90$ €; \ $24t+30=126\ \to\ t=4$ h\\[4pt]
Rectangle de perímetre 36 cm amb àrea màxima: \ $A(x)=x(18-x)$, vèrtex $x=9$: \ quadrat, 81 cm²\\[4pt]
Viatge de 240 km: \ $t(v)=\frac{240}{v}$, \ $t(80)=3$ h; \ $\frac{240}{v}=2{,}5\ \to\ v=96$ km/h\\[4pt]
2\,000 € al 3\,\% compost: \ $C(t)=2\,000\cdot1{,}03^t$, \ $C(5)\approx2\,318{,}55$ €\\[4pt]
Cinema: 500 espectadors a 6 €; per cada 0,50 € de rebaixa, 50 més: \ $I(n)=(6-0{,}5n)(500+50n)=$\\ \hspace*{1em}$=-25n^2+50n+3\,000$; \ vèrtex $n=1$: \ 5,50 €, 550 espectadors i \textbf{3\,025 €}
\end{exemple}

\vfill
\begin{truc}
Si la $x$ està en un interval $[p,q]$, el màxim o el mínim està en el \textbf{vèrtex} (si és dins) o en els \textbf{extrems}: calculem $f(p)$, $f(q)$ i $f(V_x)$ i comparem.
\end{truc}

\vfill

\vfill
% ============================================================ resum
\resumtitol
\begin{tcbraster}[raster columns=2,raster equal height=rows,raster column skip=14pt,raster row skip=26pt,raster before skip=16pt]
\rtile{7.1}{Lineals}{$y=mx+n$: recta. $m$ pendent, $(0,n)$.}
\rtile{7.2}{Punt-pendent}{$y-y_0=m(x-x_0)$. \ $m=\frac{y_2-y_1}{x_2-x_1}$.}
\rtile{7.3}{Quadràtica}{$y=ax^2$: vèrtex $(0,0)$. $a>0$ $\cup$, $a<0$ $\cap$.}
\rtile{7.4}{Paràboles}{$V_x=\frac{-b}{2a}$. Talls, vèrtex i taula.}
\rtile{7.5}{Radicals}{$\sqrt x$: $D=[0,+\infty\ro$.\\[1ex] $\sqrt{x-a}+b$ comença en $(a,b)$.}
\rtile{7.6}{Proporcionalitat inversa}{$y=\frac kx$: hipèrbola. Asímptotes, els eixos.}
\rtile{7.7}{Translacions}{$y=\frac k{x-a}+b$: asímptotes $x=a$ i $y=b$.}
\rtile{7.8}{Exponencials}{$a^x$: per $(0,1)$, $R=\lo0,+\infty\ro$.\\[1ex] $C\cdot a^t$.}
\rtile{7.9}{Problemes}{Reconeix el tipus de funció. Imatge, antiimatge o vèrtex.}
\end{tcbraster}
\end{document}
